\chapter*{\hfill{\centering  ЗАКЛЮЧЕНИЕ}\hfill}
\addcontentsline{toc}{chapter}{ЗАКЛЮЧЕНИЕ}

На основе полученных данных был сделан вывод о том, что при размерности матрицы более 9 алгоритм муравьиной колонии по времени работает более эффективно, чем алгоритм решения задачи коммивояжёра полным перебором. Однако последний упомянутый алгоритм гарантированно даёт самый эффективный результат.

Цель данной работы была выполнена. 
Для достижения поставленной цели были выполнены следующие задачи:
\begin{itemize}[label=---]
	\item описаны вышеперечисленные алгоритмы решения задачи коммивояжёра;
	\item написана программа, реализующая вышеперечисленные алгоритмы решения задачи коммивояжёра;
	\item проведена параметризация муравьиного алгоритма;
	\item проверена работоспособность программы, выполнено модульное тестирование;
	\item проведён анализ временных затрат вышеперечисленных алгоритмов.
\end{itemize}